%%%%%%%%%%%%%%%%%%%%%%%%%%%%%%%%%%%% Chapter Template

\chapter{Introdução} 	% Main chapter title
\label{Chapter1} 		% For referencing the chapter elsewhere, usage \ref{Chapter1}

%%%%%%%%%%%%%%%%%%%%%%%%%%%%%%%%%%%%

%RESUMO Feito no capítulo 0
%ABSTRACT Feito no capítulo 0
%ÍNDICE GERAL Feito
%LISTA DE SIGLAS/ABREVIATURAS (opcional) Feito
%ÍNDICE DE FIGURAS Feito
%ÍNDICE DE TABELAS Feito
%Normalmente, as páginas iniciais ou preliminares - agradecimentos, resumos, índice, etc. devem ser numeradas em algarismos romanos minúsculos; Feito
%1. Introdução • Contextualização sobre a temática em estudo;

%2. Revisão da Literatura (c/ base numa análise bibliométrica c/ mínimo de 30 referências bibliográficas) • Identificar os principais contributos e seus autores; • Identificar áreas de influência e interesse. Maior ênfase em 7 referências bibliográficas;Peso maior!
%3. BenchMarking: Case Studies • Caso(s) de Sucesso: • Proposta de Valor / Diferenciação da Concorrência; • Identificar determinantes do sucesso Menor Peso!
%4. Conclusões • Apontar as limitações do estudo; • Destacar principais contributos; • Potencialidades futuras.
%BIBLIOGRAFIA feito
%APÊNDICES / ANEXOS Feito
%Short Bio dos Membros do Grupo Feito
%%%%%

%Revisto e atualizado a 8 de Janeiro de 2026 por Eduardo Junqueira
A terminologia "Urban Logistics in Smart Cities", em inglês, ou "Logística Urbana em Cidades Inteligentes", em português, refere-se ao título deste documento  que representa a temática em estudo.
O objetivo é fazer uma análise teórica desta temática e uma análise em um exemplo de um caso prático, onde será aplicado conceitos de metodologia restritos e casos de estudo fornecidos consequentemente,  onde professor António Amaral, orientou através da plataforma académica do \gls{isep}: \gls{moodle} da \gls{uc} de \gls{logis}\cite{MoodleEnunciado}.

%Revisto e atualizado a 8 de Janeiro de 2026 por Eduardo Junqueira
\section{Contextualização}
\label{sec:contextualizacao	}

%Contextualização sobre a temática em estudo
A logística urbana é uma componente decissiva para o funcionamento eficiente das cidades mais desenvolvidas, especialmente no contexto em inglês \textit{"Smart Cities"}.
Com o crescimento populacional e a urbanização acelerada nas grandes cidades, a procura por soluções logísticas eficazes tornou-se cada vez mais crítica nos dias de hoje.

Em Portugal as \textit{"Urban Areas"}, representam as grandes cidades urbanas com maior índice populacional, onde existe um crescente aumento populacional todos os anos devido a fatores externos de imigração, de trabalho empresarial e de formação académica \cite{INE2022Censos}.

O  \gls{ine}, divulgou que os Munícipios de Portugal Continetal e Ilhas com maior número populacional registado pelos censos de 2021 foram:
Lisboa, Sintra, Vila Nova de Gaia, Porto, Cascais, Loures e Braga \cite{INE2022Censos}.

Assim, as \textit{"Smart Cities"}, que também representam as cidades urbanas mas por definição são cidades inteligentes que utilizam tecnologias avançadas  tais como por exemplo: \textit{\gls{ai}}, \textit{\gls{iot}}, \textit{\gls{bigdata}}, \textit{\gls{ml}}, \textit{\gls{dl}}, \textit{\gls{bct}} entre outras.

Desta forma, a Logística tem como objetivo implementar estas tecnologias no que diz respeito  ás \textit{"Urban Logistics"} e \textit{"Smart Logistics"}, com o papel fundamental de melhorar os problemas existentes nas áreas urbanas tornando-as mais inteligentes  e sustentáveis para um futuro próximo.

%%%%%%%%%%%%%%%%%%%%%%%%%%%%%%%%%%%%%%%%%%%%%%%%%%%%%%%%%
%Revisto e atualizado a 8 de Janeiro de 2026 por Eduardo Junqueira
\section{Metodologia}
\label{sec:metodologia}
A metodologia planeada para este estudo antes da revisão da literatura, capítulo \ref{Chapter2} , foi a seguinte:
\begin{enumerate}
    \item \textbf{\gls{query}:} ferramenta \gls{wos} com filtro de palavras-chave ou \textit{keywords} relevantes, anexo A: \ref{anexo:base_dados_resultados_query}.
    \item \textbf{\gls{vosviewer}:} análise visual dos dados obtidos, anexo B: \ref{anexo:figuras_complementares_wos}.
    \item \textbf{Seleção:} artigos mais relevantes para o tema em estudo com base em métricas específicas, anexo C: \ref{anexo:figuras_complementares_vosviewer}.
\end{enumerate}

{\setlength{\parindent}{0pt}
\textbf{ 1 - \gls{query}:}
}
A análise, deste título do documento, será feita com base em filtros específicos de pesquisa.
A utilização de operadoes lógicos como \verb|AND| e \verb|OR| foram essenciais para a construção desta \gls{query} e a identação dos parênteses foi também essencial para a organização do filtro, global, que é subdivido por vários sub-filtros, tal como é demonstrado no anexo A: \ref{anexo:base_dados_resultados_query}.

A ferramenta \gls{wos} \textit{(Web of Science)} é uma base de dados online de pesquisa científica e bibliográfica adequada neste estudo de mestrado do \gls{isep}  sendo assim a única base de dados utilizada para a revisão de literatura no capítulo: \ref{Chapter2}.
O filtro de pesquisa e filtros de pesquisa avançados, foram necessários com o uso de sinónimos e palavras dedicadas ao estudo que se assemelham a esta temática.
Uma das funcionalidades mais importantes do \gls{wos} é após o resultado da pesquisa a possibilidade de exportar os dados obtidos para um ficheiro \verb|.txt|, \gls{txt} que é compatível com a ferramenta \gls{vosviewer} e também para ficheiros \verb|.csv|, \gls{csv} que são compatíveis com o \gls{excel} para exportação bibliográfica dos dados.

A análise desta \gls{query} resultou assim em 536 artigos relevantes para o tema em estudo e a análise de resultados está presente no Anexo B: \ref{anexo:figuras_complementares_wos}.
O repositório público do \gls{github} onde a informação foi guardada e sequencialmente atualizada ao longo de todo o trabalho está disponível, assim,  esta informação e todas as seguintes da metodologia deste trabalho serão guardadas \cite{JunqueiraDevEduardoRepositório}.
Como existe uma diferença, de dois anos, existe um lapso de que no  ano passado 2025 existia 536 resultados obtidos no \gls{wos}  e que, este ano, 2026 existem apenas 533 artigos, a informação da ferramenta destes dados é o \gls{wos} que pode variar consoante a verificação destes dados futuramente. Assim, o repositório do \gls{github} serve de consulta da data até ao momento que foi obtido os dados em 2025 \cite{JunqueiraDevEduardoRepositório}.

Uma outra funcionalidade que foi essencial nesta fase de análise foi a análise dos próprios resultados obtidos no \gls{wos}.
As figuras abaixo referenciadas apresentam esta análise.

A \gls{fig} \ref{fig:wos_query_results} no Anexo B \ref{anexo:figuras_complementares_wos}, apresenta os resultados das categorias das áreas de estudo obtidas no \gls{wos} com a \gls{query} definida.
A \gls{fig} \ref{fig:wos_query_years} no Anexo B \ref{anexo:figuras_complementares_wos}, apresenta os resultados dos anos de publicação dos artigos obtidos no \gls{wos} com a \gls{query} definida.

\textbf{De notar que:} esta \gls{query} e esta análise no \gls{wos} foi implementada após várias tentaivas de filtros e palavras-chave com diferentes operadores lógicos até chegar a este resultado final que apresenta 536 artigos mais específicos e de acordo com o esperado e com o enunciado do trabalho académico \cite{MoodleEnunciado}.


{\setlength{\parindent}{0pt}
\textbf{ 2 - \gls{vosviewer}:}
}
Um  \textit{ software open-source} utilizado nesta metodologia específica para criar  e implementar o resultado obtido do filtro do \gls{wos} em um ficheiro \verb|.txt|.
Assim, é possível criar mapas de rede de dois tipos de análise: \textit{Co-authorship} e \textit{Co-occurrence}, segue o conceito de um \textit{mind-map} mas gráficamente é mais complexo e intuitivo de visualizar para base de dados maiores e mais complexas onde é quase impossível inserir ou retirar manualmente dados \cite{VOSviewer2026}.
Ao criar estes mapas é necessário primeiramente seguir um conjunto de passos que estão disponíveis na documentação do \gls{moodle} \cite{MoodleEnunciado}.

\textbf{De notar que:} todos os mapas criados são de acordo com os dados bibliográficos do ficheiro \verb|.txt|, exportado do \gls{wos} da \gls{query} definida, que se encontram também disponível no repositório do \gls{github} \cite{JunqueiraDevEduardoRepositório}.

Os mapas \textit{Co-authorship} permitem identificar relações entre autores, instituições e países.

Os mapas \textit{Co-occurrence} permitem identificar relações entre ocorrências de palavras-chave/\textit{keywords} presentes nos artigos, sendo estes são os de maior ênfase e escolhidos nesta análise.


Algumas das figuras mais ilustrativas que foram obtidas nas primeiras fases de crição destes dois mapas estão presentes no Anexo C: \ref{anexo:figuras_complementares_vosviewer}, demonstram outras tentavivas de mapas com diferentes parâmetros mais relevantes.


A \gls{fig} \ref{fig:vosviewer_final-authorship} apresenta a visualização  de tipo de análise de coautoria , teste apenas para análise de autores (não foi considerado): \textit{Software VOSviewer: Co-authorship Analysis / Authors}.

A \gls{fig} \ref{fig:vosviewer_interm-co-occurrence} apresenta a visualização de tipo de análise de coocorrência, teste intermédio: \textit{Software VOSviewer: Co-occurrence Analysis / All Keywords}.


O mapa de \textit{Co-occurrence} (análise de coocorrência, em português) da  \gls{fig} \ref{fig:vosviewer_network-final-co-occurrence}, foi o que apresentou melhor resultado, após algumas tentativas de análise no \gls{vosviewer} com a \gls{query} previamente definida, obteve o melhor resultado: \textit{Items}: 22 (palavras-chave), \textit{Links}: 166 (ligações), \textit{Total link strength}: 1638 (força total dos links) e \textit{Clusters}: 4 (grupos identificados) \cite{VosViewerResult2025}.



{\setlength{\parindent}{0pt}
\textbf{ 3 - Seleção:}
}
A seleção dos artigos definitivos, foi feita com o resultado da \gls{fig} \ref{fig:vosviewer_network-final-co-occurrence}, onde foi então validada a \gls{query} e depois foi feita a seleção dos artigos consoante os mais citados, mas do tipo \textit{Open-access} (acesso livre) na plataforma do \gls{wos}.
Foram selecionados 7 artigos mais relevantes para o tema em estudo, que serão referenciados no capítulo \ref{Chapter2}.
%%%%%%%%%%%%%%%%%%%%%%%%%%%%%%%%%%%%

%Revisto e atualizado a 8 de Janeiro de 2026 por Eduardo Junqueira
\section{Estrutura e Organização}
\label{sec:estrutura_organizacao}

Este documento tem como estrutura a seguinte abaixo descrita:

{\setlength{\parindent}{0pt}

\textbf{ Capítulo: \ref{Chapter1} - Introdução:} contextualização da temática em estudo, metodologia aplicada e por fim estrutura e organização do documento.

\textbf{ Capítulo: \ref{Chapter2} - Revisão da Literatura:} análise bibliométrica de artigos, fontes bibliográficas e autores mais relevantes.

\textbf{ Capítulo: \ref{Chapter3} - Caso de Estudo:} análise de 1 caso de estudo com maior detalhe na proposta de valor e diferenciação da concorrência, identificação dos determinantes de sucesso e de fracasso.

\textbf{ Capítulo: \ref{Chapter4} - Conclusão:} síntese dos principais pontos abordados, limitações, contributos e potencialidades futuras desta temática.

%\textbf{ Capítulo: \ref{AppendixA} - Apêndices/Anexos:} materiais complementares que suportam o conteúdo do relatório.

\textbf{ Capítulo: \ref{Chapter5} - Nota Bibliográfica:} breve biografia do estudante que realizou este trabalho.
}
No ínicio após a entrega de temas aos grupos, existiu alguns contratempos e dificuldades com o grupo. Inicialmente planeado no que diz respeito á organização de tarefas, mas levou a uma maior esforço final em Janeiro deste trabalho pelo membro que ficou sozinho Eduardo Junqueira.
Desta maneira, este trabalho  visa seguir a estrutura proposta inicialmente, mas com o ajuste de que apenas serão referidos 7 artigos dos 20 propostos iniciais, e como é de prever o número de referências citadas neste relatório será menor do que o definido.
Assim, este trabalho que foi inicialmente planeado e estruturado por 4 membros do grupo, acabou no final por ser apenas um membro a fazer o trabalho todo.
%%%%%%%%%%%%%%%%%%%%%%%%%%%%%%%%%%%%
